\section{Conclusion}

\label{sec:conclusion}

We have presented the Harness-Hacker Protocol, a framework for safe agent
self-modification organized into four trust tiers with formal isolation
guarantees, validated modification gates, and bounded session overhead. The
core technical contributions are:

\begin{enumerate}[leftmargin=1.5em]
    \item \textbf{Worktree Isolation Theorem:} Modifications in isolated git
          worktrees provably cannot affect the running system state, providing
          the foundation for all safety guarantees.
    \item \textbf{Four-Tier Taxonomy:} Runtime tools (7-day TTL), extensions
          (test-gated, persistent), skills (knowledge-only), and cross-component
          modifications (CI + human-gated) cover the full spectrum of harness
          improvement at calibrated trust levels.
    \item \textbf{Virtuous Cycle Formalization:} The feedback loop connecting
          telemetry, GRPO extraction, modification proposals, and improved
          capability is modeled as a fixed-point iteration that converges to an
          optimal tool set under mild assumptions.
    \item \textbf{Risk Analysis with Formal Bounds:} Each identified failure
          mode has a formally bounded blast radius under HHP constraints.
    \item \textbf{Empirical Validation:} Five virtuous cycles on the Hanzo Dev
          benchmark produced a $17.4\%$ improvement in task completion, with a
          $9.3$ percentage point gap between specialized and general tools
          motivating the entire endeavor.
\end{enumerate}

The broader implication is a shift in how we conceptualize the agent--harness
relationship. The harness is not a fixed environment that the agent must learn
to navigate; it is mutable infrastructure that the agent can improve, within
safety boundaries that are both formally specified and practically enforced.
Agents that can improve their own tools will, over time, outperform agents
that cannot---and HHP provides the framework to let them do so safely.

% ============================================================================
% APPENDIX
% ============================================================================
\appendix

